\section{Introduction}\label{sec:introduction}

Renewable-resource institutions operate on schedules. Surveys generate observations; assessments translate observations into advice; a rule sets a catch or effort command; and that command is held, or updated only at the next authorized review. In fisheries, the review interval may be annual or multi-year, while the ecological state evolves continuously between decisions (Butterworth and Punt, 1999; Butterworth, 2007; Punt et al., 2016). The institutional clock is therefore part of the dynamical system, not merely a reporting convention.

Two common mathematical substitutes discard that architecture. A continuous-delay equation lets feedback act continuously on a lagged signal. An annual difference equation compresses all between-review dynamics into one step. Neither is generally the map implemented by an institution that samples, decides, and holds. Sampled-data control has long treated the zero-order hold as an operator with its own spectral properties (Chen and Francis, 1995; \AA{}str\"om and Wittenmark, 1997; Ne\v{s}i\'c and Teel, 2004). Resource-governance models do not always preserve that distinction, even though stability can change when the decision operator changes.

A second problem is causal promotion: a spectral peak, cohort cycle, or crisis-driven monitoring record does not by itself identify the sign, timing, or responsiveness of a governance rule. The skeptical re-analysis of the Russell Cycle, for example, found that a long record treated as periodic did not resolve a true periodicity (McManus, Licandro, and Coombs, 2016). Here the same diagnostic discipline is applied to claims about institutional feedback. Periodicity is not treated as a mechanism label.

This paper makes four contributions. First, it specifies a sample-and-hold model in which observation, assessment, review, decision, deployment, ecological-response, and memory times are kept separate. Second, it proves forward invariance of the nonnegative state space and states precisely what the rapid-review limit does establish: finite-horizon approximation under regularity conditions, not finite-interval stability. Third, it compares an exponential frozen-assessment controller update with forward Euler on a common illustrative ecological plant and records how crossings move with the operator and parameter vector. Fourth, it re-examines the empirical layer: a multiplicity-controlled 42-stock spectral screen, a documented but coarse review-interval proxy test, a curated 32-system case search, and the northern-cod record.

The focal controller is deliberately diagnostic rather than normative. It raises extraction effort in response to a perceived shortfall, representing a mobilising or extractive feedback. A protective comparator reduces effort when the perceived deficit rises. The numerical coefficients are inherited illustrative conventions, not field estimates; the exercise therefore identifies operator dependence within a specified model class, not a calibrated review interval for a real fishery. This distinction matters especially because the baseline crossing lies close to the unit circle and moves substantially under small parameter perturbations.

The paper proceeds as follows. Section~\ref{sec:methods} defines the model and evidence procedures. Section~\ref{sec:results} reports the two mathematical properties, the operator-specific crossing record, sensitivity, and empirical results. Section~\ref{sec:discussion} states what transfers to management design and what remains unidentified. Detailed proofs, case records, data tables, and computational audits are in the Supplementary Information.

\begin{figure}[htbp]
\centering
\includegraphics[width=0.98\linewidth]{fig_concept_v38.png}
\caption{Three representations of institutional timing. Continuous-delay feedback acts on a lagged signal; an annual-step model compresses the between-decision dynamics; sample-and-hold governance observes and revises at review times and holds the command between them. Panel B is schematic: the computed markers identify the annual spectral radii and the exponential-update crossing, not a universal institutional response curve.}
\label{fig:concept}
\end{figure}

